\documentclass{ximera}
\title{Limits and Derivatives, Part 2, Trigonometric Functions}


\newcommand{\pskip}{\vskip 0.1 in}

\begin{document}
\begin{abstract}
Derivatives of some trig functions.
\end{abstract}
\maketitle


\pskip

\section{The Sine Function}

\begin{example}  \label{Ex:9u4erdfg9}
\begin{enumerate}
\item Zoom in near the origin on the graph of the function
\[
   y = f(\theta) = \sin\theta \, , \, -2\pi \leq \theta \leq 2\pi.
\]
below to approximate the value of the derivative
\[
     \frac{dy}{d\theta}\Big|_{\theta = 0}  =   \frac{d}{d\theta} \left( \sin \theta \right)\Big|_{\theta = 0} .
\]
Geometrically, the derivate is the slope of the tangent line to the curve $y=\sin\theta$ at the origin.


\begin{onlineOnly}
    \begin{center}
\desmos{tordzgkr1za}{900}{600}
\end{center}
\end{onlineOnly}

Desmos activity available at
\href{https://www.desmos.com/calculator/ordzgkr1za}{151: Sine Graph Basic}

\item Find a simplifed expression for the function $m=r(\theta_1)$ that gives the average rate of change of the function $f(\theta)$ over the interval between $\theta = 0$ and $\theta = \theta_1$. State the domain of this function.

\item Sketch by hand a graph of the function $m=r(\theta_1)$. Drag point $B$ in the graph above to help.

\item Input your expression for the function $m$ in Line 1 of the worksheet. Then open and activate the average rate of change folder in Line 2. Describe what you see? What does this suggest about the value of the derivative
\[
   \frac{d}{d\theta} \left( \sin \theta \right)\Big|_{\theta = 0} = \lim_{\theta_0 \to 0} \frac{\sin\theta_0}{\theta_0} ?
\]

\end{enumerate}

\end{example}


\begin{example} \label{Exldf3r3ffs}
Let
\[
  g(\theta) = \sin \left( \frac{1}{3} \theta \right) \, , \, -6\pi \leq \theta \leq 6\pi.
\]

\begin{enumerate}

\item Start by graphing this function by hand. Describe the transformation that takes the graph of $y=f(\theta) = \sin\theta$ to the graph of 
\[
 y = g(\theta) = f\left(  \frac{1}{3}\theta \right).
\]

\item Use part (a) to guess the value of the derivative
\[
    \frac{d}{d\theta} \left( \sin \left(\frac{1}{3}\theta\right) \right)\Big|_{\theta = 0} .
\]
Explain your thinking.

\item Find a simplifed expression for the function $m=r(\theta_1)$ that gives the average rate of change of the function $g(\theta)$ over the interval between $\theta = 0$ and $\theta = \theta_1$. State the domain of this function.

\item Use the fact that
\[
 \lim_{\theta\to 0} \frac{\sin\theta}{\theta} = 1
\]
and the algebra of limits to evaluate the derivative in part (b).

\end{enumerate}
\end{example}



\begin{example} \label{ExMf90923111}

Repeat parts (a), (b), and (d) of Example 1 for the derivative
\[
    \frac{d}{d\theta} \left( \sin \theta \right)\Big|_{\theta = \pi/6}.
\]
Use the graph of the function $y=\sin\theta$ below to help. 


\begin{onlineOnly}
    \begin{center}
\desmos{bf2ic534ik}{900}{600}
\end{center}
\end{onlineOnly}

Desmos activity available at
\href{https://www.desmos.com/calculator/bf2ic534ik}{151: Sine Graph Basic 2}

\end{example}



\begin{example}
Let 
\[
   f(\theta) = \sin\theta .
\]

\begin{enumerate}
\item Find an expression for the average rate of change of the function $f(\theta)$ between angles $\theta = \theta_1$ and $\theta = \theta_2$.

\item Use part (a), the algebra of limits, and the fact that
\[
  \lim_{\theta\to 0} \frac{\sin\theta}{\theta} = 1
\]
to find an expression for the derivative
\[
   \frac{d}{d\theta} \left( \sin \theta \right)\Big|_{\theta = \theta_1} .
\]

You'll need the identity
\[
   \sin B - \sin A = 2\sin\left( \frac{1}{2} \left( B+A  \right)\right) \cos\left( \frac{1}{2} \left( B+A  \right)\right) .
\]

\item Use the result of part (c) to evaluate the derivative
\[
 \frac{d}{d\theta} \left( \sin \theta \right)\Big|_{\theta =\pi/6 } .
\]
Compare the exact value with your approximation from Example 3.
\end{enumerate}

\end{example}


\section{A Building's Shadow}

\begin{example} \label{ExLfm3209}
The function
\[
   s = f(\theta) = 50\cot \theta \, , \, 0<\theta\leq \pi/2 ,
\]
expresses the length of the shadow of a $50$-foot tall building in terms of the sun's inclination angle.

\begin{enumerate}
\item Find  an expression for the average rate of change of the function $f(\theta)$ between angles $\theta = \theta_1$ and $\theta = \theta_2$.

\item  Use part (a), the algebra of limits, and the fact that
\[
  \lim_{\theta\to 0} \frac{\sin\theta}{\theta} = 1
\]
to find an expression for the derivative
\[
    \frac{ds}{d\theta} \Big|_{\theta = \theta_1} =   \frac{d}{d\theta} \left( 50 \cot \theta \right)\Big|_{\theta = \theta_1} .
\]

You'll need the identity
\[
   \sin(B-A) = \sin B \cos A - \cos B \sin A.
\]
\end{enumerate}

\item Express the derivative in part (b) in terms of the length
\[
  s = f(\theta_1) 
\]
of the shadow.

\item Use part (c) to find the exact value of the derivative
\[
   \frac{ds}{d\theta} \Big|_{s=100} .
\]
Compare with your work in the chapter

\href{https://ximera.osu.edu/calcone/Calculus1/Shadow1/Shadow1}{Shadow of a Building, Part 1}

\end{example}





\end{document}